% ============================================================
%  QUESTIONS D'ENTRETIEN — 30 Questions + Reponses
%  Module : Authentification et Gestion des Roles
%  Etudiant : Ghoulam Mohamed Said
%  Compatible Overleaf — pdfLaTeX
% ============================================================

\documentclass[12pt,a4paper,oneside]{report}

% ── Encodage & langue ──────────────────────────────────────
\usepackage[utf8]{inputenc}
\usepackage[T1]{fontenc}
\usepackage[french]{babel}
\usepackage{lmodern}

% ── Mise en page ───────────────────────────────────────────
\usepackage[top=2.5cm,bottom=2.5cm,left=3cm,right=2.5cm]{geometry}
\usepackage{setspace}
\onehalfspacing
\usepackage{parskip}
\setlength{\parindent}{0pt}

% ── Couleurs ───────────────────────────────────────────────
\usepackage{xcolor}
\definecolor{GoldenGlow}{RGB}{212,175,55}
\definecolor{SoftEspresso}{RGB}{62,39,35}
\definecolor{AccentGray}{RGB}{120,110,100}
\definecolor{EasyGreen}{RGB}{34,139,34}
\definecolor{MedBlue}{RGB}{30,80,160}
\definecolor{HardRed}{RGB}{180,30,30}

% ── Tableaux ───────────────────────────────────────────────
\usepackage{booktabs}
\usepackage{tabularx}
\usepackage{array}

% ── Boites colorees ────────────────────────────────────────
\usepackage[most]{tcolorbox}
\tcbuselibrary{breakable,skins}

\newtcolorbox{qbox}[2]{
  colback=#1!6,
  colframe=#1,
  fonttitle=\bfseries\color{white},
  coltitle=#1!30!black,
  title=#2,
  arc=4pt,
  boxrule=0.8pt,
  left=6pt,right=6pt,top=4pt,bottom=4pt,
  breakable
}

\newtcolorbox{abox}{
  colback=gray!5,
  colframe=AccentGray!60,
  arc=3pt,
  boxrule=0.5pt,
  left=6pt,right=6pt,top=4pt,bottom=4pt,
  breakable
}

% ── Sections ───────────────────────────────────────────────
\usepackage{titlesec}
\titleformat{\chapter}[display]
  {\normalfont\huge\bfseries\color{SoftEspresso}}
  {\color{GoldenGlow}\chaptertitlename\ \thechapter}
  {10pt}{\Huge}
  [\vspace{4pt}\color{GoldenGlow}\titlerule[1.5pt]]

\titleformat{\section}
  {\normalfont\Large\bfseries\color{SoftEspresso}}
  {\color{GoldenGlow}\thesection}{1em}{}

% ── En-tetes ───────────────────────────────────────────────
\usepackage{fancyhdr}
\pagestyle{fancy}
\fancyhf{}
\fancyhead[L]{\small\color{SoftEspresso}\nouppercase{\leftmark}}
\fancyhead[R]{\small\color{GoldenGlow}Questions d'Entretien}
\fancyfoot[C]{\small\color{AccentGray}\thepage}
\renewcommand{\headrulewidth}{0.4pt}

% ── Utilitaires ────────────────────────────────────────────
\usepackage{enumitem}
\usepackage{amsmath}
\usepackage{microtype}
\usepackage{hyperref}
\hypersetup{colorlinks=true,linkcolor=SoftEspresso}

% ── Code inline ────────────────────────────────────────────
\usepackage{listings}
\lstset{
  basicstyle=\ttfamily\small,
  backgroundcolor=\color{gray!8},
  frame=single, framerule=0.4pt,
  rulecolor=\color{GoldenGlow},
  breaklines=true,
  showspaces=false,
  numbers=none,
  literate=
    {é}{{\'{e}}}1 {è}{{\`{e}}}1 {ê}{{\^{e}}}1
    {à}{{\`{a}}}1 {â}{{\^{a}}}1 {ù}{{\`{u}}}1
    {î}{{\^{\i}}}1 {ô}{{\^{o}}}1 {ç}{{\c{c}}}1
}

% ============================================================
\begin{document}
% ============================================================

% ════════════════════════════════════════════════════════════
%  PARTIE I — QUESTIONS
% ════════════════════════════════════════════════════════════

\chapter*{Questions d'Entretien}
\addcontentsline{toc}{chapter}{Questions d'Entretien}
\markboth{Questions d'Entretien}{Questions d'Entretien}

\begin{tcolorbox}[colback=GoldenGlow!8,colframe=GoldenGlow,
  arc=4pt,boxrule=0.7pt]
\textbf{Module~:} Authentification et Gestion des Rôles ---
SGUI Ibn Khaldoun\\
\textbf{Étudiant~:} Ghoulam Mohamed Said\\
\textbf{Structure~:} 10 questions faciles · 10 intermédiaires ·
10 difficiles\\
\textbf{Réponses~:} Partie II (page~\pageref{ch:reponses})
\end{tcolorbox}

\vspace{0.5cm}

% ─────────────────────────────────────
\section*{\textcolor{EasyGreen}{Niveau 1 --- Facile}}
% ─────────────────────────────────────

\begin{enumerate}[
  leftmargin=1.5cm,
  label=\textcolor{EasyGreen}{\textbf{Q\arabic*.}},
  itemsep=10pt
]

\item Qu'est-ce que l'\textbf{authentification}~? En quoi diffère-t-elle
      de l'\textbf{autorisation}~?

\item Qu'est-ce qu'un \textbf{JSON Web Token (JWT)}~? De quelles parties
      est-il composé~?

\item Pourquoi utilise-t-on \textbf{bcrypt} pour hacher les mots de passe
      plutôt que MD5 ou SHA-1~?

\item Qu'est-ce que le \textbf{RBAC} (Role-Based Access Control)~?
      Donnez un exemple concret tiré de votre projet.

\item Que signifie l'en-tête HTTP
      \texttt{Authorization: Bearer <token>}~? Où est-il utilisé dans
      votre application~?

\item Qu'est-ce qu'une \textbf{API RESTful}~? Citez ses trois principaux
      principes.

\item Que se passe-t-il lorsqu'un JWT expire~? Comment le frontend
      réagit-il dans votre système~?

\item Qu'est-ce que le \textbf{CORS} et pourquoi est-il important pour
      la sécurité des API~?

\item Quelle est la différence entre stocker un token dans le
      \textbf{localStorage} et dans un cookie \texttt{HttpOnly}~?

\item Qu'est-ce qu'un \textbf{middleware} dans Express.js~? Donnez un
      exemple concret tiré de votre projet.

\end{enumerate}

\vspace{0.8cm}

% ─────────────────────────────────────
\section*{\textcolor{MedBlue}{Niveau 2 --- Intermédiaire}}
% ─────────────────────────────────────

\begin{enumerate}[
  leftmargin=1.5cm,
  label=\textcolor{MedBlue}{\textbf{Q\arabic*.}},
  start=11,
  itemsep=10pt
]

\item Expliquez la \textbf{structure détaillée} d'un JWT~: que contient
      chaque partie et comment la signature garantit-elle l'intégrité~?

\item Décrivez le \textbf{flux d'authentification complet} dans votre
      système, depuis la saisie des identifiants jusqu'à l'accès à une
      ressource protégée. Citez chaque étape.

\item Qu'est-ce qu'un \textbf{ORM}~? Quels avantages \textbf{Prisma}
      offre-t-il par rapport à des requêtes SQL brutes~?

\item Comment avez-vous implémenté le \textbf{contrôle d'accès par rôle}
      dans votre API~? Montrez un exemple avec \texttt{requireRole}.

\item Qu'est-ce que l'\textbf{injection SQL}~? Comment Prisma vous
      protège-t-il contre cette attaque~?

\item Qu'est-ce que le \textbf{rate-limiting}~? Dans quel contexte
      l'avez-vous utilisé et pourquoi est-ce important~?

\item Qu'est-ce que le \textbf{Context API} de React~? Comment
      l'avez-vous utilisé pour gérer l'état d'authentification~?

\item Comment fonctionne un \textbf{intercepteur Axios}~? Quel problème
      double résout-il dans votre application~?

\item Qu'est-ce qu'une \textbf{SPA} (Single Page Application)~? Quelles
      implications cela a-t-il pour la gestion de l'authentification~?

\item Expliquez la différence entre un \textbf{Access Token} et un
      \textbf{Refresh Token}. Avez-vous implémenté les deux dans votre
      projet~?

\end{enumerate}

\vspace{0.8cm}

% ─────────────────────────────────────
\section*{\textcolor{HardRed}{Niveau 3 --- Difficile}}
% ─────────────────────────────────────

\begin{enumerate}[
  leftmargin=1.5cm,
  label=\textcolor{HardRed}{\textbf{Q\arabic*.}},
  start=21,
  itemsep=10pt
]

\item Quelles sont les \textbf{failles de sécurité} liées au stockage du
      JWT dans le localStorage~? Quelle alternative auriez-vous pu choisir
      et pourquoi~?

\item Expliquez le problème de \textbf{révocation de token JWT}. Comment
      l'avez-vous adressé ou contourné dans votre système~?

\item Comparez l'architecture \textbf{JWT stateless} à l'architecture
      \textbf{session-based stateful}~: avantages, inconvénients, et dans
      quels cas préférer l'une ou l'autre~?

\item Qu'est-ce qu'une attaque \textbf{CSRF}~? Pourquoi les JWT en
      en-tête \texttt{Authorization} sont-ils naturellement résistants à
      cette attaque~?

\item Si un administrateur modifie le \textbf{rôle d'un utilisateur} dont
      le token est encore valide, comment votre système garantit-il que ce
      changement est pris en compte immédiatement~?

\item Expliquez le \textbf{principe du moindre privilège}
      (\textit{Least Privilege}). Comment l'avez-vous appliqué dans la
      conception de votre matrice RBAC~?

\item Comment concevriez-vous un système de \textbf{journalisation
      d'audit} pour tracer toutes les actions sensibles (connexion,
      changement de rôle, suppression de données)~?

\item Décrivez une attaque par \textbf{force brute} sur votre endpoint
      \texttt{/login}. Quelles mesures avez-vous mises en place~? Sont-elles
      suffisantes selon vous~?

\item Votre système utilise \textbf{bcrypt avec un coût de~12}. Comment ce
      facteur affecte-t-il la sécurité et les performances~? Quel coût
      recommanderiez-vous en production en~2025~?

\item Votre application doit supporter \textbf{10~000 utilisateurs
      simultanés}. Quels goulots d'étranglement anticipez-vous dans votre
      architecture actuelle et comment les résoudriez-vous~?

\end{enumerate}

\newpage

% ════════════════════════════════════════════════════════════
%  PARTIE II — REPONSES
% ════════════════════════════════════════════════════════════

\chapter*{Réponses aux Questions d'Entretien}
\label{ch:reponses}
\addcontentsline{toc}{chapter}{Reponses aux Questions d'Entretien}
\markboth{Reponses}{Reponses}

% ─────────────────────────────────────
\section*{\textcolor{EasyGreen}{Réponses --- Niveau 1 (Facile)}}
% ─────────────────────────────────────

\begin{qbox}{EasyGreen}{R1 --- Authentification vs. Autorisation}
L'\textbf{authentification} vérifie l'identité~: «~Êtes-vous bien qui
vous prétendez être~?~». Elle repose sur des identifiants (e-mail + mot
de passe) vérifiés contre la base de données.

L'\textbf{autorisation} intervient \textit{après} et détermine ce que
l'utilisateur peut faire~: «~Avez-vous le droit d'accéder à cette
ressource~?~». Un enseignant authentifié n'est pas autorisé à accéder
à la liste de tous les utilisateurs (route admin uniquement).
\end{qbox}

\begin{qbox}{EasyGreen}{R2 --- JWT et ses trois parties}
Un JWT (RFC 7519) est composé de trois parties séparées par un point~:
\begin{enumerate}
  \item \textbf{Header}~: algorithme de signature
        (\texttt{\{alg:HS256, typ:JWT\}}) en Base64Url.
  \item \textbf{Payload}~: claims~: \texttt{sub} (identifiant),
        \texttt{role}, \texttt{iat} (émission), \texttt{exp}
        (expiration). Encodé en Base64Url mais \textbf{non chiffré}.
  \item \textbf{Signature}~: HMAC-SHA256(header + payload, SECRET).
        Garantit l'intégrité du token.
\end{enumerate}
\end{qbox}

\begin{qbox}{EasyGreen}{R3 --- bcrypt vs MD5/SHA-1}
MD5 et SHA-1 sont rapides~: un GPU peut tester des milliards de
combinaisons par seconde.

\textbf{bcrypt} est conçu pour les mots de passe~:
\begin{itemize}
  \item \textbf{Sel aléatoire} intégré~: empêche les tables arc-en-ciel.
  \item \textbf{Facteur de coût 12}~: environ 250~ms par vérification,
        rendant la force brute prohibitivement longue.
\end{itemize}
\end{qbox}

\begin{qbox}{EasyGreen}{R4 --- RBAC}
Le RBAC assigne des permissions à des \textit{rôles}, et les utilisateurs
héritent des permissions de leur rôle.

Dans notre SGUI~:
\begin{itemize}
  \item \texttt{admin}   --- accès complet à toutes les routes.
  \item \texttt{teacher} --- propose des sujets PFE, gère son jury.
  \item \texttt{student} --- soumet des réclamations, consulte les
        annonces.
\end{itemize}
\end{qbox}

\begin{qbox}{EasyGreen}{R5 --- Authorization Bearer}
C'est la méthode standard (RFC 6750) pour transmettre un token d'accès.
\textit{Bearer} signifie «~porteur~»~: quiconque possède ce token est
présumé autorisé. Notre middleware \texttt{requireAuth} extrait la valeur
après «~Bearer~» et la vérifie via \texttt{jwt.verify()}.
\end{qbox}

\begin{qbox}{EasyGreen}{R6 --- API RESTful}
Une API RESTful respecte~:
\begin{itemize}
  \item \textbf{Stateless}~: chaque requête contient toutes les
        informations nécessaires.
  \item \textbf{Méthodes HTTP standard}~: GET, POST, PUT, DELETE.
  \item \textbf{Ressources identifiées}~: URI unique par ressource
        (\texttt{/api/users/:id}).
\end{itemize}
\end{qbox}

\begin{qbox}{EasyGreen}{R7 --- JWT expiré}
Le backend retourne \texttt{401 TOKEN\_INVALID\_OR\_EXPIRED}.
L'intercepteur Axios détecte ce code et~:
\begin{enumerate}
  \item Supprime le token du \texttt{localStorage}.
  \item Redirige automatiquement vers \texttt{/login}.
\end{enumerate}
\end{qbox}

\begin{qbox}{EasyGreen}{R8 --- CORS}
Le CORS bloque par défaut les requêtes HTTP vers un domaine différent.
Sans configuration, le frontend sur \texttt{localhost:5173} ne pourrait
pas appeler l'API sur \texttt{localhost:5000}. En production, CORS
protège l'API contre les appels depuis des sites malveillants non
autorisés.
\end{qbox}

\begin{qbox}{EasyGreen}{R9 --- localStorage vs cookie HttpOnly}
\begin{tabular}{@{}lll@{}}
\toprule
\textbf{Critère} & \textbf{localStorage} & \textbf{Cookie HttpOnly} \\
\midrule
Accès JavaScript & Oui   & Non    \\
Risque XSS       & Elevé & Faible \\
Risque CSRF      & Nul   & Modéré \\
Simplicité       & Haute & Moyenne\\
\bottomrule
\end{tabular}

\vspace{4pt}
Nous avons choisi le localStorage pour la simplicité dans ce contexte
académique. En production, un cookie \texttt{HttpOnly + Secure +
SameSite=Strict} est recommandé.
\end{qbox}

\begin{qbox}{EasyGreen}{R10 --- Middleware Express}
Un middleware s'intercale dans le pipeline de traitement d'une requête,
avec accès à \texttt{req}, \texttt{res} et \texttt{next()}.

Notre \texttt{requireAuth}~:
\begin{itemize}
  \item Vérifie la présence du token JWT.
  \item Si valide~: attache \texttt{req.user} et appelle
        \texttt{next()}.
  \item Si invalide~: retourne \texttt{401} et stoppe le pipeline.
\end{itemize}
\end{qbox}

\vspace{0.5cm}

% ─────────────────────────────────────
\section*{\textcolor{MedBlue}{Réponses --- Niveau 2 (Intermédiaire)}}
% ─────────────────────────────────────

\begin{qbox}{MedBlue}{R11 --- Structure détaillée d'un JWT}
\begin{itemize}
  \item \textbf{Header}~: \texttt{\{alg:HS256, typ:JWT\}} en Base64Url.
        Précise l'algorithme.
  \item \textbf{Payload}~: \texttt{\{sub:uuid, role:admin, iat:1234,
        exp:1235\}}. \textit{Attention~: non chiffré, seulement encodé.}
        Ne jamais y mettre de données sensibles.
  \item \textbf{Signature}~: HMAC-SHA256(base64url(header) + "." +
        base64url(payload), SECRET). Toute modification invalide la
        signature, détectée et rejetée par le serveur.
\end{itemize}
\end{qbox}

\begin{qbox}{MedBlue}{R12 --- Flux d'authentification complet}
\begin{enumerate}
  \item L'utilisateur saisit e-mail + mot de passe sur LoginPage.
  \item \texttt{POST /api/auth/login} envoyé au backend.
  \item \texttt{prisma.user.findUnique(\{email\})} recherche le compte.
  \item \texttt{bcrypt.compare(password, passwordHash)} vérifie le mdp.
  \item \texttt{jwt.sign(\{sub, role\}, SECRET, \{expiresIn:'24h'\})}
        génère le token.
  \item Le token est stocké dans le \texttt{localStorage}.
  \item Chaque requête suivante~: l'intercepteur Axios ajoute
        \texttt{Authorization: Bearer <token>}.
  \item \texttt{requireAuth} decode et vérifie le token.
  \item \texttt{requireRole} vérifie que le rôle est autorisé.
  \item Le contrôleur traite la requête et retourne la réponse.
\end{enumerate}
\end{qbox}

\begin{qbox}{MedBlue}{R13 --- ORM et avantages de Prisma}
Un \textbf{ORM} fait le pont entre le code objet et la base relationnelle.

Avantages de Prisma~:
\begin{itemize}
  \item \textbf{Schéma déclaratif}~: \texttt{schema.prisma} comme
        source de vérité unique.
  \item \textbf{Client typé}~: autocomplétion, détection d'erreurs à la
        compilation.
  \item \textbf{Requêtes paramétrées}~: protection native contre
        l'injection SQL.
  \item \textbf{Migrations versionnées}~: évolution contrôlée du schéma.
\end{itemize}
\end{qbox}

\begin{qbox}{MedBlue}{R14 --- Implémentation RBAC avec requireRole}
\texttt{requireRole} est une fonction d'ordre supérieur~:
\texttt{requireRole(['admin'])} retourne un middleware qui vérifie
\texttt{req.user.role}.

\begin{lstlisting}
router.get('/users',
  requireAuth,            // Verifier le JWT
  requireRole(['admin']), // Verifier le role
  listUsers               // Executer le controleur
);
\end{lstlisting}
Si le rôle ne correspond pas~: \texttt{403 FORBIDDEN}.
\end{qbox}

\begin{qbox}{MedBlue}{R15 --- Injection SQL et protection Prisma}
L'injection SQL consiste à insérer du code SQL malveillant dans un champ
(ex.~: \texttt{' OR '1'='1} dans le champ email pour bypasser
l'authentification).

Prisma protège via des \textbf{requêtes préparées}~: les données
utilisateur ne sont jamais concaténées dans une chaîne SQL --- elles
sont toujours passées comme paramètres typés.
\end{qbox}

\begin{qbox}{MedBlue}{R16 --- Rate-limiting}
Le rate-limiting limite le nombre de requêtes d'un client dans un
intervalle donné. Nous l'appliquons sur \texttt{POST /auth/login}~:
maximum 20 tentatives par IP par 15 minutes.

Au-delà~: \texttt{429 Too Many Requests}. Cela rend les attaques par
dictionnaire prohibitivement lentes.
\end{qbox}

\begin{qbox}{MedBlue}{R17 --- Context API React}
Le Context API permet de partager des données sans \textit{prop
drilling}. Notre \texttt{AuthContext} expose~:
\begin{itemize}
  \item \texttt{user}~: objet utilisateur ou \texttt{null}.
  \item \texttt{login(email, password)}~: authentifie et stocke le token.
  \item \texttt{logout()}~: efface le token et réinitialise l'état.
  \item \texttt{loading}~: indique si la vérification initiale est en
        cours.
\end{itemize}
Tout composant dans \texttt{AuthProvider} accède via \texttt{useAuth()}.
\end{qbox}

\begin{qbox}{MedBlue}{R18 --- Intercepteur Axios}
Un intercepteur s'exécute automatiquement sur chaque requête/réponse~:
\begin{itemize}
  \item \textbf{Request}~: lit le token et l'injecte dans l'en-tête
        \texttt{Authorization} --- évite de le répéter dans chaque
        appel API.
  \item \textbf{Response}~: si \texttt{401}, efface le token et
        redirige vers \texttt{/login} automatiquement.
\end{itemize}
\end{qbox}

\begin{qbox}{MedBlue}{R19 --- SPA et authentification}
Dans une SPA, il n'y a pas de rechargement entre navigations, donc~:
\begin{itemize}
  \item L'état d'auth est géré côté client (Context API).
  \item Les routes privées sont bloquées par \texttt{ProtectedRoute}.
  \item \textbf{Important}~: le backend doit \textit{toujours} vérifier
        le token --- un attaquant peut appeler l'API directement sans
        passer par React.
\end{itemize}
\end{qbox}

\begin{qbox}{MedBlue}{R20 --- Access Token vs. Refresh Token}
\begin{itemize}
  \item \textbf{Access Token}~: durée courte (15 min--24h), utilisé sur
        chaque requête. En cas de vol, exposition limitée.
  \item \textbf{Refresh Token}~: durée longue (7--30 jours), utilisé
        uniquement pour renouveler l'Access Token sans reconnexion.
\end{itemize}
Notre système utilise uniquement un Access Token de 24h (pour simplifier
l'implémentation). Les Refresh Tokens sont une perspective
d'amélioration.
\end{qbox}

\vspace{0.5cm}

% ─────────────────────────────────────
\section*{\textcolor{HardRed}{Réponses --- Niveau 3 (Difficile)}}
% ─────────────────────────────────────

\begin{qbox}{HardRed}{R21 --- Vulnerabilites du localStorage}
Le localStorage est accessible depuis JavaScript~: une attaque XSS peut
lire et exfiltrer le token.

\textbf{Alternative}~: cookie avec~:
\begin{itemize}
  \item \texttt{HttpOnly}~: inaccessible depuis JavaScript.
  \item \texttt{Secure}~: transmis uniquement en HTTPS.
  \item \texttt{SameSite=Strict}~: protection anti-CSRF.
\end{itemize}
Choix localStorage justifié par la simplicité dans ce contexte
académique.
\end{qbox}

\begin{qbox}{HardRed}{R22 --- Revocation de token JWT}
Un JWT est valide jusqu'à expiration même après déconnexion --- problème
inhérent au caractère stateless.

\textbf{Stratégies~:}
\begin{enumerate}
  \item \textbf{Durée courte}~: 15--60 min + Refresh Tokens.
  \item \textbf{Blacklist Redis}~: tokens révoqués stockés en cache.
        Réintroduit de l'état côté serveur.
  \item \textbf{Vérification BDD}~: sur routes critiques, vérifier que
        \texttt{user.isActive} et que le rôle n'a pas changé.
\end{enumerate}
Notre système applique la stratégie 3 sur les opérations sensibles.
\end{qbox}

\begin{qbox}{HardRed}{R23 --- JWT Stateless vs. Session Stateful}
\begin{tabular}{@{}lll@{}}
\toprule
\textbf{Critere}   & \textbf{JWT Stateless} & \textbf{Session Stateful} \\
\midrule
Etat serveur       & Aucun                  & BDD / Redis requis        \\
Scalabilite        & Horizontale naturelle  & Session partagee requise  \\
Revocation         & Difficile              & Triviale                  \\
Performance        & Verification locale    & Lecture BDD par requete   \\
Payload visible    & Oui (Base64)           & Non (opaque)              \\
\bottomrule
\end{tabular}

\vspace{4pt}
JWT est préféré pour les microservices. Les sessions conviennent mieux
quand la révocation instantanée est critique (ex.~: banques).
\end{qbox}

\begin{qbox}{HardRed}{R24 --- CSRF et resistance naturelle du JWT Bearer}
CSRF incite le navigateur à envoyer une requête authentifiée à l'insu
de l'utilisateur (fonctionne avec les cookies, envoyés automatiquement).

Les JWT en en-tête \texttt{Authorization} résistent car~:
\begin{enumerate}
  \item Le navigateur n'envoie \textbf{jamais automatiquement} un
        en-tête personnalisé en cross-origin.
  \item Seul le JavaScript légitime peut lire le localStorage.
  \item Un site malveillant ne peut pas accéder au localStorage d'un
        autre domaine (Same-Origin Policy).
\end{enumerate}
\end{qbox}

\begin{qbox}{HardRed}{R25 --- Changement de role et token en cours}
Sur les routes critiques, une vérification en base de données confirme
que le rôle est toujours valide~:

\begin{lstlisting}
const dbUser = await prisma.user.findUnique({
  where: { id: req.user.sub }
});
if (!dbUser?.isActive || dbUser.role !== req.user.role)
  return res.status(403).json({ error: 'ROLE_CHANGED_RELOGIN' });
\end{lstlisting}

Cela a un coût en performance mais garantit la cohérence sur les
opérations sensibles.
\end{qbox}

\begin{qbox}{HardRed}{R26 --- Principe du moindre privilege}
Chaque entité ne doit disposer que des permissions strictement
nécessaires.

Application dans le SGUI~:
\begin{itemize}
  \item Un \texttt{student} accède uniquement à ses propres réclamations.
  \item Un \texttt{teacher} modifie uniquement ses sujets PFE (et
        seulement quand \texttt{status === 'propose'}).
  \item Même l'\texttt{admin} ne retourne jamais le
        \texttt{passwordHash} dans les réponses API.
\end{itemize}
\end{qbox}

\begin{qbox}{HardRed}{R27 --- Journalisation d'audit}
Conception recommandée~:
\begin{enumerate}
  \item \textbf{Modele}~:
        \texttt{(id, userId, action, targetId, metadata, timestamp)}.
  \item \textbf{Actions}~: connexion/déconnexion, changement de rôle,
        création/suppression de compte, validation PFE, décision
        réclamation.
  \item \textbf{Stockage}~: table PostgreSQL en append-only ou service
        externe (Elasticsearch).
  \item \textbf{Middleware}~: appelé dans chaque contrôleur sensible.
\end{enumerate}

\begin{lstlisting}
await prisma.auditLog.create({
  data: {
    userId,
    action,
    metadata: JSON.stringify(metadata),
    timestamp: new Date()
  }
});
\end{lstlisting}
\end{qbox}

\begin{qbox}{HardRed}{R28 --- Attaque par force brute sur /login}
\textbf{Deroulement}~: l'attaquant envoie des milliers de requêtes
\texttt{POST /login} en testant des mots de passe pour un email cible.

\textbf{Mesures implementees}~:
\begin{itemize}
  \item Rate-limiting~: 20 requêtes / 15 min par IP.
  \item bcrypt coût 12~: \textasciitilde250 ms par vérification.
  \item Messages génériques~: \texttt{INVALID\_CREDENTIALS} sans
        révéler si l'email existe.
\end{itemize}

\textbf{Insuffisances}~: le rate-limiting par IP est contournable via
des proxies. Améliorations~: CAPTCHA après N échecs, blocage temporaire
du compte, notification par email.
\end{qbox}

\begin{qbox}{HardRed}{R29 --- Facteur de cout bcrypt en 2025}
Le facteur de coût détermine $2^{\text{coût}}$ itérations~:

\begin{tabular}{@{}lll@{}}
\toprule
\textbf{Cout} & \textbf{Duree} & \textbf{Recommandation} \\
\midrule
10 & \textasciitilde64 ms  & Trop rapide, a eviter  \\
12 & \textasciitilde250 ms & Recommande en 2025     \\
14 & \textasciitilde1 s    & Systemes tres sensibles\\
\bottomrule
\end{tabular}

\vspace{4pt}
Règle~: choisir le coût tel que le hachage prenne 100--500 ms sur le
serveur de production, à réévaluer tous les 2 ans.
\end{qbox}

\begin{qbox}{HardRed}{R30 --- Scalabilite a 10 000 utilisateurs}
\textbf{Goulots anticipes}~:
\begin{enumerate}
  \item \textbf{Base de donnees}~: PostgreSQL single-instance devient
        bottleneck. Solution~: PgBouncer (connection pooling), replicas
        de lecture, index optimises.
  \item \textbf{bcrypt sur /login}~: 250 ms $\times$ 10~000 = CPU
        significatif. Solution~: worker threads ou file asynchrone
        (Bull + Redis).
  \item \textbf{JWT stateless}~: verification locale, pas de probleme.
  \item \textbf{API}~: scalable horizontalement derriere un load
        balancer (Nginx) devant N instances Node.js.
\end{enumerate}

\textbf{Architecture cible}~:
Nginx $\rightarrow$ N $\times$ Node.js $\rightarrow$ PostgreSQL +
PgBouncer + Redis.
\end{qbox}

% ============================================================
\end{document}
% ============================================================
